% Canonical native-TeX manuscript; compile from the repository root.
% Publication layout changes are reversible in the import manifest.
% Original Subset Sum paper style, recovered without changing its typography.
% Original main.tex SHA-256:
% ff77ed5a42d72a79197b6881a35cf7982457667b21d88bf5a1eeaea61689945a
% Original subset-sum.pdf SHA-256:
% bdeb95400d450f64cfc27ebbae255718b9a6a26607d0c1375091a4c7b3a4a431
\documentclass[11pt,a4paper]{article}
\usepackage[T1]{fontenc}
\usepackage[utf8]{inputenc}
\usepackage{lmodern}
\usepackage[a4paper,margin=25mm]{geometry}
\usepackage{amsmath,amssymb,amsthm}
\usepackage{booktabs,longtable,array}
\usepackage{microtype}
\usepackage[numbers,sort&compress]{natbib}
\usepackage[unicode,colorlinks=true,linkcolor=blue,citecolor=blue,urlcolor=blue,hypertexnames=false]{hyperref}
\urlstyle{same}
\setlength{\emergencystretch}{2em}
\setlength{\parskip}{3pt}
\clubpenalty=10000
\widowpenalty=10000
\displaywidowpenalty=10000
\allowdisplaybreaks

\usepackage{mathtools}
\newtheorem{theorem}{Theorem}[section]
\newtheorem{lemma}[theorem]{Lemma}
\newtheorem{corollary}[theorem]{Corollary}
\newcommand{\F}{\mathbb F}
\newcommand{\SV}{\operatorname{SV}}
\newcommand{\wt}{\operatorname{wt}}
\newcommand{\pdeg}{\operatorname{pdeg}}
\title{Direction Grids for All-Characteristic\\
Affine-Orbit Identity Testing}
\author{Yilin Liu \and GPT-6 Astra Pro}
\date{27 October 2026}
\begin{document}\maketitle

\begin{abstract}
We refine the finite grids in the companion all-rank affine-orbit
hitting theorem. At each rank stage, it suffices to preserve one minor
for each of at most $(n+2)(n+3)/2$ distinct contiguous intervals.
Digit polarization of a maximum-digit-weight leading monomial gives
a nonzero tensor that is separately Frobenius-linear in every direction.
Each interval test therefore excludes at most $n-1$ values per direction,
independently of all degrees. The resulting bound is
$Q^{K}(nD+1)(D+1)^K$, where
$Q=(n-1)(n+2)(n+3)/2+1$ and $K=K(w)=2^{O(w^2)}$.
Choosing globally between this argument and ordinary interpolation
replaces $Q$ by $\min\{Q,(n-1)D+1\}$ and restores the working-field
requirement $|E|\ge nD+1$.
\end{abstract}

\noindent\textbf{Status and dependency.}
The mathematical plaintext and this formal addendum have separate
model-assisted internal reconstructions and exact-source audits.
Their mathematical content is frozen after those internal reviews.
This is not external peer review or a publication-priority claim.
The companion manuscript supplies the all-characteristic symbolic
rank-descent theorem and its Artin and Hasse lemmas. This addendum proves
the finite-grid refinement and does not replace those structural proofs.
The canonical companion is \texttt{manuscripts/general-affine.tex}
in this problem directory. Its research-session source before the
documented publication layout changes has SHA-256
\begin{center}\footnotesize\ttfamily
9864136e134bbc9c3f35d8d0f1fb7595cf40c397e01d0247173956795f0c695f.
\end{center}

\section{Model and quantitative statements}

Let $n,w,D\ge1$, let $\F$ be any field, and let the nonzero oracle
polynomial have a representation
\begin{equation}\label{eq:model}
 f(X)=uA_1(\ell_1(X)+c_1)\cdots A_n(\ell_n(X)+c_n)v,
 \qquad \deg f\le D.
\end{equation}
The univariate matrices $A_i$ are $w$-by-$w$ over $\F$, with arbitrary
entry degrees and generic ranks. The forms $\ell_i$ are the coordinates
of an invertible linear transformation. The order, affine transformation,
and representation are unknown. The degree promise concerns only the
scalar polynomial.

For $1\le r<w$ put
\begin{align}
 W_r&=\binom wr,\qquad
 h_r=(r+2)\prod_{s=r+2}^w
                   \left(\binom{s-1}{r-1}+1\right)-2,\label{eq:h}\\
 k_r&=(2h_r+2)(W_r^2-1),\qquad
 K=K(w)=\sum_{r=1}^{w-1}k_r.\label{eq:K}
\end{align}
Set $K(1)=0$. Empty products are one. The companion proves
$K(w)=2^{O(w^2)}$ and that
\[
 f\bigl(\gamma(t)+\SV_K(z,y)\bigr)\ne0,
 \qquad \gamma_i(t)=t^i,
\]
in every characteristic. Here $\lambda_1,\ldots,\lambda_n$ are the
Lagrange polynomials for $n$ distinct nodes in a working field $E$
extending $\F$, and
\[
 G_i(t,z,y)=t^i+\sum_{j=1}^K z_j\lambda_i(y_j).
\]

Define
\begin{equation}\label{eq:Q}
 M=\frac{(n+2)(n+3)}2,\qquad Q=(n-1)M+1.
\end{equation}
These integers depend only on $n$, and $Q\le2n^3$.

\begin{theorem}[Degree-independent direction grids]\label{thm:grid}
Let $\mathcal Y,\mathcal T,\mathcal Z\subseteq E$ be arbitrary sets
of respective sizes $Q,nD+1,D+1$. The evaluations
\[
 \left\{G(t,z,y):
 t\in\mathcal T,\quad z_j\in\mathcal Z,\quad
 y_j\in\mathcal Y\ (1\le j\le K)\right\}
\]
hit every nonzero polynomial \eqref{eq:model}. Their number is at most
\begin{equation}\label{eq:count}
 Q^K(nD+1)(D+1)^K.
\end{equation}
It suffices that $|E|\ge\max\{Q,nD+1\}$. The sets may overlap.
For fixed $w$, the count is $O_w(n^{3K+1}D^{K+1})$.
\end{theorem}

The result also holds with the initial curve
\begin{equation}\label{eq:eta}
 \eta_i(t)=t\lambda_i(t)
\end{equation}
in place of $\gamma_i(t)$. The corresponding evaluations lie in the
image of pure $\SV_{K+1}$, with its first pair specialized to
$(z_0,y_0)=(t,t)$.

\begin{corollary}[Smaller field and the better of two grids]\label{cor:min}
Put
\[
 B=\min\{Q,(n-1)D+1\}.
\]
Over a working field of size at least $nD+1$, a Cartesian hitting
grid for either initial curve uses at most
\begin{equation}\label{eq:mincount}
 (nD+1)\bigl((D+1)B\bigr)^K
\end{equation}
queries. Every direction uses the same size-$B$ set: the choice
between the two arguments is made globally.
\end{corollary}

Extension-field evaluations require the corresponding black-box access.
The statements do not promise that base-field-valued queries alone
always suffice. Nonzero constants need one query. If $n=1$, ordinary
univariate interpolation gives the sharper elementary count $D+1$,
independently of the displayed general upper bounds.

\section{A nonzero Frobenius direction tensor}

For characteristic $p>0$, let $\wt_p(a)$ be the sum of the base-$p$
digits of $a$. The digit weight of $X^\alpha$ is
$\sum_i\wt_p(\alpha_i)$. In characteristic zero use ordinary degree,
so every digit copy below has size one.

\begin{lemma}\label{lem:tensor}
Let $K_0$ be a coefficient field containing $E$, and let $H(X)$ be a
nonzero ordinary homogeneous polynomial of degree $m>0$ and digit
degree at most $k$. There are $k_0\le k$ powers
$q_1,\ldots,q_{k_0}$ of the characteristic, or ones in characteristic
zero, with $\sum_jq_j=m$, such that the coefficient of
$\prod_jz_j^{q_j}$ in
\[
 H\left(\sum_{j=1}^{k_0}z_jV_j\right)
\]
is a nonzero tensor
\begin{equation}\label{eq:tensor}
 T(V_1,\ldots,V_{k_0})
   =\sum_{i_1,\ldots,i_{k_0}}b_{i_1,\ldots,i_{k_0}}
                        \prod_{j=1}^{k_0}V_{j,i_j}^{q_j}.
\end{equation}
After substituting $V_{j,i}=\lambda_i(y_j)$, the tensor remains
nonzero. For each direction $y_j$, at most $n-1$ values in $E$
can make it identically zero in all the other directions.
\end{lemma}
\begin{proof}
Choose a nonzero monomial $X^\alpha$ of \emph{maximum} digit weight
$k_0$ among the monomials of $H$. List its digit copies as the
$q_j$, with repetition according to each digit. Their sum is $m$.

In the Frobenius expansion of a source monomial, every one of the
$k_0$ positive amplitude exponents in the selected coefficient must
receive at least one digit copy. No source monomial has more than
$k_0$ copies. Thus every contributor has exactly $k_0$ copies, with
exactly one assigned to each amplitude. In particular, no carry
within an amplitude produces a mixed direction monomial. This proves
the separate Frobenius-linearity in \eqref{eq:tensor}.

Assign the chosen monomial's copies to their corresponding direction
coordinates. The resulting direction-monomial coefficient is the
nonzero coefficient of $X^\alpha$ times
\[
 \prod_{i,e}\alpha_{i,e}!,
 \qquad \alpha_i=\sum_e\alpha_{i,e}p^e
\]
in positive characteristic. Every digit is less than $p$, so all
these factorials are nonzero. In characteristic zero the coefficient
is instead $\prod_i\alpha_i!$, also nonzero. No division by $m!$
is used. A different source monomial cannot yield this same direction
monomial, since its exponents sum back to $\alpha$. Hence $T\ne0$.

The functions $\lambda_i(y_j)^{q_j}$ are linearly independent over
$K_0$. Their coefficient matrix on the powers $y_j^{a q_j}$ is
the $q_j$-th Frobenius power of the invertible Lagrange coefficient
matrix. In characteristic zero it is that matrix itself. Applying
these injective linear maps in each tensor factor preserves $T\ne0$.

View the resulting tensor as a polynomial in $y_j$ over the rational
function field in all the other directions. It is a nonzero linear
combination of $\lambda_i(y_j)^{q_j}$. Over an algebraic closure of
that coefficient field it is a $q_j$-th power of a polynomial of
degree at most $n-1$. It therefore has at most $n-1$ distinct roots.
This proves the assertion about values making the entire remaining
tensor zero. Root extraction is only a proof device; neither
perfection nor an algorithm for taking these roots is assumed.
\end{proof}

The maximum-digit-weight choice is necessary for this argument.
Selecting an arbitrary lower-weight monomial would allow additional
digit copies from other monomials to carry inside one amplitude and
would not establish the special tensor form \eqref{eq:tensor}.

\begin{lemma}[A finite family of tests]\label{lem:family}
Suppose $M_0$ nonzero scalar tests over $K_0$ have first homogeneous
components of digit degree at most $k$. Then one can specialize all
$k$ fresh SV directions to any prescribed set of
$(n-1)M_0+1$ distinct elements of $E$ so that every test remains
nonzero as a function of the old parameters and the fresh amplitudes.
\end{lemma}
\begin{proof}
A test with nonzero constant term imposes no condition. For every
other test choose the nonzero tensor of Lemma~\ref{lem:tensor},
using at most $k$ directions; additional directions are ignored.
Specialize directions one at a time. As long as a tensor is nonzero,
it excludes at most $n-1$ choices for the next direction. The union
of the forbidden choices has size at most $(n-1)M_0$, so one of the
prescribed values keeps all tensors nonzero. Repeat.

The chosen tensor is an amplitude coefficient of the original test.
Higher ordinary homogeneous terms have larger total amplitude degree
and cannot change it. Thus each test survives. This is an existential
proof for an explicit Cartesian grid; the algorithm does not identify
tests, tensors, or the good specialization.
\end{proof}

\section{Counting the rank-descent conditions}

We recall exactly the structural facts from the companion proof that
are used here. It adjoins two constant scalar-endpoint factors to the
$n$ original matrices. At rank stage $r$, every augmented interval
of generic rank greater than $r$ already retains its rank at the
current symbolic shift. A generic rank-$r$ target interval is cut into
shortest rank-$r$ blocks and a higher-rank tail:
\[
 I=B_1\cdots B_qT,\qquad B_j=C_jA_{i_j}.
\]
Each $C_j$ and $T$ is empty or has every interval of rank greater
than $r$. The required rank-$r$ tests are
\begin{equation}\label{eq:tests}
 B_j,\qquad B_jB_{j+1}\ (j<q),\qquad B_qT.
\end{equation}
Each is a contiguous interval of the same augmented word.

The companion's Artin-box and filtered local-Hasse arguments prove
that \emph{every nonzero $r$-minor} of any word in \eqref{eq:tests}
has its entire first homogeneous residual component of digit degree
at most $k_r$ from \eqref{eq:K}. This is the filtered conclusion that
permits Lemma~\ref{lem:family}; a low-weight witness somewhere in the
series would not suffice.

There are $N=n+2$ factors in the augmented word, hence at most
$N(N+1)/2=M$ nonempty contiguous intervals. For every generic
rank-$r$ target take its required intervals \eqref{eq:tests}, and
deduplicate the entire list by interval. Choose \emph{one} nonzero
$r$-minor of each distinct required interval.

Preserving that one minor certifies rank at least $r$; rank at most
$r$ follows from the identically vanishing larger minors. The chosen
minor works in every target decomposition where the interval occurs.
For its leading-component bound, choose any one such decomposition.
There is therefore no extra factor for the number of minors, the
number of target occurrences, or repeated decompositions.

At most $M$ scalar conditions need to be preserved at each stage.
Lemma~\ref{lem:family} applies with the same grid size
$Q=(n-1)M+1$ for every one of the stage's $k_r$ directions.

\section{Fixing directions while retaining the rank invariant}

This specialization must be done within the symbolic descent, not
merely after a rational proof has been frozen. At the start of a
stage, fix every earlier direction to its previously selected element
of $E$, while keeping $t$ and all earlier amplitudes algebraically
independent. The current coefficient field is
\[
 K_0=E(t,\text{earlier amplitudes}).
\]
Every latent shift has the form
\[
 C_i=c_i+\ell_i(t,t^2,\ldots,t^n)
               +\text{a polynomial in earlier amplitudes}.
\]
Its $t$-part is nonconstant, since $\ell_i$ is a nonzero linear form
and $t,t^2,\ldots,t^n$ are independent polynomials over $E$.
Thus each $C_i$ remains transcendental over the original coefficient
field and over $E$, in every characteristic. This is the precise
local-Hasse premise; no derivative condition is needed.

By induction all already prepared higher-rank interval minors survive
as functions of the formal old parameters. The common direction grid
preserves every new rank-$r$ test as a function of these parameters
and the new amplitudes. Old nonzero minors remain nonzero because
setting new amplitudes to zero recovers them. The companion's
adjacent-block chain argument consequently preserves every rank-$r$
target after the directions have been fixed.

All unit pivots, coefficient-space bases, and constant cutwise gauges
may be chosen anew over the current coefficient field. They are proof
devices and introduce no additional direction conditions. In particular,
we do not require an earlier rational expression for a gauge to
retain its denominator under a later specialization. The actual tests
are minors of original polynomial matrix intervals. Their analysis
uses coefficients in the current old-parameter field, not new-direction
denominators.

Iterating through the ranks gives at least one full choice of all $K$
directions from the common set $\mathcal Y$ for which $f(G)$ is a
nonzero polynomial in $t,z_1,\ldots,z_K$. Its individual degrees are
at most $nD$ in $t$ and $D$ in each amplitude. The final
$\mathcal T,\mathcal Z$ grids therefore hit it. This proves
Theorem~\ref{thm:grid}. Auxiliary representation or minor degrees do
not enter this last interpolation step.

\section{The alternative curve and the minimum-size rule}

For \eqref{eq:eta}, every nonzero latent linear combination is
$c+tq(t)$, with $q$ a nonzero combination of the independent Lagrange
polynomials. It is nonconstant in every characteristic, and $\eta$
has degree at most $n$. The entire stagewise argument and final
degree counts are unchanged. This proves the pure-map specialization
stated after Theorem~\ref{thm:grid}.

\begin{proof}[Proof of Corollary~\ref{cor:min}]
If $Q\le(n-1)D+1$, use Theorem~\ref{thm:grid} for all directions.
Then $Q\le nD+1$, so its field requirement is simply $|E|\ge nD+1$.
If $Q>(n-1)D+1$, use instead ordinary Cartesian interpolation on
the full nonzero substituted scalar polynomial from the companion
theorem. Every direction has degree at most $(n-1)D$, every amplitude
degree at most $D$, and the initial parameter degree at most $nD$.
This gives the stated size with $B=(n-1)D+1$.

The choice is global: one uses the same size-$B$ direction grid for
all pairs, justified entirely by one of these two arguments. No
unproved mixture of stagewise degree and root-count bounds is needed.
Both cases yield \eqref{eq:mincount} over the advertised field.
\end{proof}

For fixed width, the degree-independent branch has exponent $K+1$
in the scalar degree instead of the basic $2K+1$ bound. Its dependence
on $n$ is larger. No optimality is asserted. The affine transformation
must remain invertible; singular projections and unrestricted repeated
variables are outside the model. All mathematical conclusions retain
the internal-review status of the companion proof.

\end{document}
